% Generated draft: requires mathematical and source review.
\documentclass[11pt]{article}
\usepackage[margin=25mm]{geometry}
\usepackage{fontspec}
\setmainfont{Latin Modern Roman}
\usepackage{amsmath,amssymb,mathtools,unicode-math}
\setmathfont{Latin Modern Math}
\usepackage{enumitem,needspace,xurl,hyperref}
\setlength{\parindent}{0pt}
\setlength{\parskip}{5pt}
\setlength{\emergencystretch}{3em}
\newcommand{\N}{\mathbb{N}}
\newcommand{\Z}{\mathbb{Z}}
\newcommand{\Q}{\mathbb{Q}}
\newcommand{\R}{\mathbb{R}}
\newcommand{\C}{\mathbb{C}}
\newcommand{\F}{\mathbb{F}}
\newcommand{\A}{\mathbb{A}}
\newcommand{\PP}{\mathbb{P}}
\newcommand{\E}{\mathbb{E}}
\newcommand{\eps}{\varepsilon}
\newcommand{\dd}{\,\mathrm{d}}
\begin{document}
\section*{UnsolvedMath ID: 800028}
\textbf{Source problem code:} catalogue entry
\subsection{Definitions and mathematical statement}
A semigroup is a nonempty set with an associative multiplication. A zero is an element \(0\) satisfying \(0s=s0=0\) for every element \(s\). Use the standard convention that a semigroup \(S\) with zero is \(0\)-simple if \(S\ne\{0\}\), \(S^2\ne\{0\}\), and its only nonempty two-sided ideals are \(\{0\}\) and \(S\). For a nonzero idempotent \(e\in S\), its maximal subgroup is the group of units of the monoid \(eSe\), with identity \(e\).

A pseudovariety of groups is a nonempty class \(\mathcal V\) of finite groups closed under isomorphisms, subgroups, homomorphic images, and finite direct products. In the usual group language \(\{\cdot,{}^{-1},1\}\), define
\[
\operatorname{Th}_{\forall}(\mathcal V)
=\left\{\forall x_1\cdots\forall x_n\,\varphi:
\varphi\text{ is quantifier-free and }
G\models\forall x_1\cdots\forall x_n\,\varphi
\text{ for every }G\in\mathcal V\right\}.
\]
Decidability of this universal theory means that an algorithm determines, for every universal sentence in this language, whether it belongs to \(\operatorname{Th}_{\forall}(\mathcal V)\). The source does not specify a logical signature; this definition makes the usual group-language convention explicit.

For a finite group \(G\) and a finite nonempty set \(I\), the Brandt semigroup \(B(G,I)\) has underlying set \((I\times G\times I)\cup\{0\}\), with zero absorbing and multiplication
\[
(i,g,j)(k,h,\ell)=
\begin{cases}
(i,gh,\ell),&j=k,\\
0,&j\ne k.
\end{cases}
\]
Here the convention allows \(|I|=1\); the supplied statement gives no restriction on \(|I|\).

Under the standard interpretation of a finite \(0\)-simple semigroup ``over'' a group as one whose nonzero maximal subgroups are isomorphic to that group, put
\[
\begin{aligned}
\mathcal C_{\mathcal V}
&=\{S:S\text{ is a finite }0\text{-simple semigroup and every nonzero maximal subgroup of }S\text{ belongs to }\mathcal V\},\\
\mathcal B_{\mathcal V}
&=\{B(G,I):G\in\mathcal V,\ 1\le |I|<\infty\},
\end{aligned}
\]
where the second class is understood up to isomorphism.

\textbf{Source questions, attributed to S. I. Kublanovsky.} For a pseudovariety \(\mathcal V\) of groups whose universal theory is decidable:
\begin{enumerate}
\item Is \(\mathcal C_{\mathcal V}\) completely solvable?
\item Is \(\mathcal B_{\mathcal V}\) completely solvable?
\end{enumerate}
The phrase ``completely solvable'' is retained verbatim as an unresolved technical qualification. The supplied record does not define it or identify the decision problems it entails, so these questions cannot yet be converted into fully specified algorithmic assertions. The source also does not specify whether algorithms may depend on the fixed pseudovariety \(\mathcal V\), or whether a uniform procedure from an effective description of \(\mathcal V\) and its universal-theory decision procedure is intended.
\subsection{Short English statement}
S. I. Kublanovsky asks whether finite \(0\)-simple semigroups, and separately finite Brandt semigroups, are ``completely solvable'' when their underlying groups lie in a pseudovariety with decidable universal theory. The supplied record leaves ``completely solvable'' undefined, so the intended algorithmic conclusion requires clarification.
\subsection{Sources}
\begin{itemize}
\item Local UnsolvedMath registry: permanent ID \(800028\), source code \texttt{catalogue entry}, title \emph{The Lyapin's Notebook --- Problem 29}; supplied upstream revision \texttt{d371ee917ac9647fbb06cd6fc529deec6f3dda82}. The registry labels this record open; that label is supplied metadata, not an independently verified mathematical or literature assessment.
\item Bershadsky, Kublanovsky, and Mashevitzky, \emph{The Lyapin's Notebook: Unsolved Problems in Semigroup Theory} (2026), Problem 29, \url{https://arxiv.org/abs/2604.04763}. The problem statement specifically attributes the questions to S. I. Kublanovsky; the registry's broader proposer field names the three collection authors and gives 2026.
\item Supplied provenance records access on 2026-10-16 and extraction by \texttt{arxiv\_latex}. This access date and provenance are supplied curation information; the source has not been independently visited or verified for this draft.
\end{itemize}
\end{document}
